\begin{afterword}
\summaryhead

The post asks what a weak chess player predicts by believing that Kasparov is better than a lesser grandmaster, when the weak player's guesses about either man's moves would be the same. The answer is the final position. The belief predicts that the game will end in the class of positions that are wins for Kasparov, and the strength of the belief is the probability placed on that class. The prediction rests on knowing Kasparov's goal and his power to reach it.

A second example makes the same point. A friend driving the author to the airport through unknown streets makes turns the author cannot predict, yet the author can confidently predict arrival, from any starting point, by knowing the friend's goal, skill and knowledge of the city. Unlike step-by-step simulation or a closed-form solution, this model predicts the outcome and none of the steps. The author calls such knowledge the subject of a life's work, extends it to any ``optimization process,'' natural selection among them, and hopes to gain precise knowledge of this kind about Friendly AI.

\responsehead

The post's answer is right. A belief about an agent's competence has empirical content, and the content lies in outcomes, not in moves. Chess ratings measure skill in just this way: a rating difference predicts an expected score and nothing about the moves.

The observation is old, and the post does not say so. William James made the airport point in 1890, with iron filings and a magnet set against Romeo and Juliet: for the filings the path is fixed, for the lover ``it is the end which is fixed, the path may be modified indefinitely.'' Daniel Dennett's ``intentional stance,'' set out from 1971, predicts a system by treating it as an agent with goals, and Dennett's own summary of it takes chess programs as its example.

The puzzle that opens the post needs a caveat that the post leaves out. If the author's probabilities for each next move were the same for Kasparov and for Mr.~G in every position, those probabilities would fix the chances of every whole game, and the results could not favour Kasparov, apart from the colour each plays. So the confident outcome prediction must use knowledge that the move guesses do not. The author said this the day before, in ``Expected Creative Surprises'': the situation is possible because ``I am not logically omniscient.'' The caveat bears on the chess comparison, not on the airport. There, uncertainty about each turn is consistent with confidence about the destination. In the chess case, identical move predictions for two players cannot coherently give different predictions of the result.

The widening at the end is announced more than shown. ``Optimization process'' takes in natural selection, but the example gives only what cannot be predicted, the form of the next organism, and not the outcome that can. The post closes on the author's aims for Friendly AI.

In short: a sound account of how belief in a capable agent predicts outcomes but not steps, a point William James had made in 1890, opened with a puzzle that depends on a caveat the author stated the day before and left out here.
\end{afterword}
